\documentclass[letterpaper]{report}
\usepackage{booktabs}
\begin{document}
\begin{table}

\caption{\label{tbl-menaces}Menaces perçues sur l'environnement et
responsables identifiés (\% de Oui)}

\centering{

\fontsize{12.0pt}{14.0pt}\selectfont
\begin{tabular*}{\linewidth}{@{\extracolsep{\fill}}lrr}
\toprule
 & {\bfseries {Alaotra}} & {\bfseries {Marovoay}} \\ 
\midrule\addlinespace[2.5pt]
\multicolumn{3}{l}{{\bfseries {Perception globale}}} \\[2.5pt] 
\midrule\addlinespace[2.5pt]
Considèrent que l'environnement est menacé & 82 % & 79 % \\ 
\midrule\addlinespace[2.5pt]
\multicolumn{3}{l}{{\bfseries {Principales menaces (parmi ceux qui déclarent menacé)}}} \\[2.5pt] 
\midrule\addlinespace[2.5pt]
Feux & 84 % & 95 % \\ 
Défrichage & 56 % & 52 % \\ 
Prélèvement (faune) & 40 % & 20 % \\ 
Prélèvement (flore) & 40 % & 21 % \\ 
Prospection minière & 29 % & 6 % \\ 
Nomadisme (passage non-résidents) & 16 % & 16 % \\ 
Zébu / pâturage & 20 % & 10 % \\ 
Tourisme & 2 % & 3 % \\ 
\midrule\addlinespace[2.5pt]
\multicolumn{3}{l}{{\bfseries {Principaux responsables}}} \\[2.5pt] 
\midrule\addlinespace[2.5pt]
L'ensemble des riverains & 87 % & 70 % \\ 
Les riverains migrants & 47 % & 37 % \\ 
L'État (Fanjakana) & 10 % & 8 % \\ 
Les braconniers & 30 % & 10 % \\ 
Le gestionnaire & 9 % & 3 % \\ 
Les nomades ou semi-nomades & 22 % & 17 % \\ 
Les ONG ou entreprises & 3 % & 4 % \\ 
Les chercheurs & 0 % & 2 % \\ 
Les touristes & 1 % & 2 % \\ 
\bottomrule
\end{tabular*}
\begin{minipage}{\linewidth}
\textbf{Source : Enquête auprès des OR 2025}\\
\end{minipage}

}

\end{table}%
\end{document}
